\section*{Capítulo \ref{ch:b3:rna-regulation} --- RNAs não codificantes e regulação pós-transcricional}

\begin{solution}{exo:b3:rna-regulation:1}
miRNA: codificado no genoma como um grampo, \qty{22}{nt}, cortado pela
\omterm{def:b3:rna-regulation:mirna}{Drosha} e depois pela \omterm{def:b3:rna-regulation:rnai}{Dicer}, pareia de modo imperfeito por sua semente e
reprime a tradução e a estabilidade de muitos mensageiros. \omterm{def:b3:rna-regulation:ncrna}{siRNA}: vem
de RNA longo de fita dupla (viral, transgênico, experimental),
\qty{21}{nt}, dependente da \omterm{def:b3:rna-regulation:rnai}{Dicer}, pareia perfeitamente e dirige a
\omterm{def:b3:rna-regulation:rnai}{Argonauta} para clivar o alvo. \omterm{def:b3:rna-regulation:ncrna}{piRNA}: vem de transcritos de fita simples
de \omterm{def:b3:rna-regulation:pirna}{agrupamentos de piRNA}, \qtyrange{26}{31}{nt}, independente da \omterm{def:b3:rna-regulation:rnai}{Dicer},
carregado em proteínas PIWI na linhagem germinativa, cliva transcritos
de transpósons (pingue-pongue) e dirige o silenciamento da \omterm{def:b3:chromatin-epigenetics:nucleosome}{cromatina} do
DNA dos transpósons.
\end{solution}

\begin{solution}{exo:b3:rna-regulation:2}
A Pol II transcreve o pri-miRNA; a \omterm{def:b3:rna-regulation:mirna}{Drosha} (com a DGCR8) recorta o
grampo no núcleo; a exportina-5 exporta o pré-miRNA; a \omterm{def:b3:rna-regulation:rnai}{Dicer} corta a
alça, deixando um dúplex de \qty{22}{nt}; uma das fitas é carregada na
\omterm{def:b3:rna-regulation:rnai}{Argonauta}; a semente (nucleotídeos 2--8) pareia com um sítio na 3$'$
UTR; a \omterm{def:b3:rna-regulation:rnai}{Argonauta} recruta a maquinaria de desadenilação e de remoção do
capuz e bloqueia a iniciação --- o mensageiro é menos traduzido e decai
mais depressa.
\end{solution}

\begin{solution}{exo:b3:rna-regulation:3}
O RNA feito in vitro por uma polimerase de fago continha contaminantes
de fita dupla (a polimerase também copia a outra fita do molde e
ultrapassa as extremidades), de modo que as preparações ``senso''
carregavam o verdadeiro gatilho. Fire e Mello purificaram cada fita,
mostraram que cada uma isolada fazia pouco e depois as anelaram de
propósito: o RNA de fita dupla era ao menos cem vezes mais potente, e
algumas moléculas por célula bastavam.
\end{solution}

\begin{solution}{exo:b3:rna-regulation:4}
Privada de ferro: a IRP liga os IREs. Ferritina: a iniciação da
tradução é bloqueada e nenhuma proteína de armazenamento é feita
(controle da tradução). Receptor de transferrina: o mensageiro é
protegido de sua nuclease e se acumula, de modo que mais receptor é
feito e mais ferro é importado (controle da estabilidade do
mensageiro).
\end{solution}

\begin{solution}{exo:b3:rna-regulation:5}
$\gamma = \ln 2/\qty{120}{min} = \qty{0.0058}{min^{-1}}$; $m^{*} =
5/0.0058 \approx 870$ moléculas. Com o decaimento dobrado:
$m^{*} \approx 430$, meia-resposta $\qty{120}{min}/2 = \qty{60}{min}$.
\end{solution}

\begin{solution}{exo:b3:rna-regulation:6}
Sequência total de 3$'$ UTR $2\times 10^{7}$ nt; correspondências ao
acaso $2\times
10^{7}/4^{7} = 2\times 10^{7}/\num{16384} \approx \num{1200}$ sítios. Um alvo real se
distingue pela conservação do sítio entre espécies, por um contexto
favorável (uma vizinhança rica em AU, um sítio afastado do códon de
parada, pareamento adicional na extremidade 3$'$ do miRNA), pela
coexpressão de miRNA e alvo na mesma célula, e pelo experimento: o
mensageiro cai quando o miRNA está presente e deixa de cair quando o
sítio é mutado.
\end{solution}

\begin{solution}{exo:b3:rna-regulation:7}
(a) XX, \emph{tra} nulo: a SXL é feita, mas a TRA não, e o \emph{dsx}
toma o splicing masculino padrão --- um macho somático (estéril). (b)
XY com TRA constitutiva: a TRA-2 está presente nos dois sexos, de modo
que a DSX é feita na forma feminina --- desenvolvimento somático
feminino (uma ``pseudofêmea'' estéril, com a linhagem germinativa
seguindo outros sinais). (c) Sem \emph{dsx}: nenhuma das formas de DSX,
de modo que nenhum dos programas de diferenciação terminal é imposto
--- um intersexo de caracteres mistos ou incompletos.
\end{solution}

\begin{solution}{exo:b3:rna-regulation:8}
No estado estacionário o mensageiro, e portanto a proteína, é
proporcional à meia-vida do mensageiro:
$\qty{180}{min}/\qty{15}{min} = 12$ vezes mais proteína. A proteína é
um sinal de crescimento normal que deveria ser transitório, um pulso
após cada estímulo; feita doze vezes mais e de forma contínua, ela
impulsiona a proliferação sem estímulo. A dose e o tempo, e não a
sequência, é que a tornam oncogênica.
\end{solution}

\begin{solution}{exo:b3:rna-regulation:9}
(a) Mãe de laboratório, pai selvagem: os óvulos não carregam \omterm{def:b3:rna-regulation:ncrna}{piRNAs}
contra elementos P (a linhagem da mãe nunca os encontrou), e os
elementos P paternos transpõem livremente na linhagem germinativa da
prole --- estéril (disgênica). (b) Mãe selvagem: seus óvulos contêm
\omterm{def:b3:rna-regulation:ncrna}{piRNAs} de elemento P, que silenciam os elementos na prole --- fértil. A
assimetria é a deposição materna de \omterm{def:b3:rna-regulation:ncrna}{piRNAs}, uma herança de RNA, e não
de genes.
\end{solution}

\begin{solution}{exo:b3:rna-regulation:10}
A proteína integra seu mensageiro ao longo de sua própria vida
$\tau_{p}$. Os mensageiros vivem $\tau_{m} = 1/(\gamma + \gamma' R)$;
durante $\tau_{p}$ a proteína vê, portanto, cerca de
$\tau_{p}/\tau_{m}$ tempos de vida independentes de mensageiro, cada um
um evento aleatório de nascimento e morte. A flutuação relativa de uma
média sobre $N$ eventos independentes cai com $1/\sqrt{N}$, de modo que
o $\tau_{m}$ mais curto sob controle de \omterm{def:b3:rna-regulation:ncrna}{microRNA} --- com o mesmo
mensageiro médio, o que exige um $k$ proporcionalmente maior --- dá uma
proteína mais lisa. A célula paga pelo \omterm{def:b3:rna-regulation:ncrna}{microRNA} e pela transcrição
extra para comprar precisão e velocidade: um nível fixado por um
equilíbrio de síntese rápida e decaimento rápido é ao mesmo tempo menos
ruidoso e mais pronto a se reajustar do que o mesmo nível fixado por
síntese lenta e decaimento lento.
\end{solution}

\begin{solution}{exo:b3:rna-regulation:11}
Verme: a RNA-polimerase dependente de RNA faz novos \omterm{def:b3:rna-regulation:ncrna}{siRNAs} a partir do
próprio mensageiro-alvo, de modo que o sinal é renovado enquanto o alvo
for transcrito, sobrevive à diluição pelas divisões e, com um canal de
RNA de fita dupla entre as células, espalha-se a outros tecidos e à
linhagem germinativa por algumas gerações. Mamífero: sem amplificação,
os \omterm{def:b3:rna-regulation:ncrna}{siRNAs} são consumidos e diluídos a cada divisão, agem apenas nas
células que os receberam e se desvanecem em dias nas células em
divisão. Consequência: um fármaco de \omterm{def:b3:rna-regulation:ncrna}{siRNA} precisa ser entregue a cada
célula-alvo, funciona melhor em células que não se dividem
(hepatócitos, daí os fármacos hepáticos) e exige doses repetidas.
\end{solution}

\begin{solution}{exo:b3:rna-regulation:12}
(a) RNA funcional: destruir o transcrito depois de feito (\omterm{def:b3:rna-regulation:ncrna}{siRNA},
oligonucleotídeo antissenso) dá um fenótipo, e expressar o RNA a partir
de um transgene em outro lugar do genoma o resgata. (b) Transcrição
funcional: inserir um sinal precoce de poliadenilação que interrompa a
transcrição (ou deletar o promotor) dá o fenótipo, ao passo que
destruir o RNA pós-transcricionalmente não dá, e um transgene ectópico
não resgata --- o ato de transcrever (ou o elemento de DNA) é que
importa, como em muitos RNAs associados a intensificadores. (c) Ruído:
nenhuma das manipulações muda coisa alguma, o RNA está a menos de uma
cópia por célula, e sua sequência e sua expressão não são conservadas.
\end{solution}

\begin{solution}{pb:b3:rna-regulation:1}
\textbf{1.} $\gamma = \ln 2/40 = \qty{0.0173}{min^{-1}}$; $m^{*} =
3/0.0173 \approx 170$ mensageiros.
\textbf{2.} $\delta_{p} = \ln 2/180 = \qty{0.00385}{min^{-1}}$; $P^{*} =
\beta m^{*}/\delta_{p} = 2\times 173/0.00385 \approx 9.0\times 10^{4}$
proteínas.
\textbf{3.} Decaimento $4\gamma = \qty{0.069}{min^{-1}}$; $m^{*} = 3/0.069
\approx 43$; constante de tempo $1/0.069 = \qty{14}{min}$.
\textbf{4.} $\beta = 1$: $P^{*} = 43/0.00385 \approx 1.1\times 10^{4}$;
repressão de $8$ vezes ($4$ pelo decaimento, $2$ pela tradução).
\textbf{5.} Meia-resposta do mensageiro $\ln 2/(4\gamma) = \qty{10}{min}$.
A proteína decai então com sua própria meia-vida, \qty{3}{h}: é a
degradação da proteína que limita o interruptor.
\textbf{6.} Sim: o mensageiro cai apenas a $1/4$ (ainda facilmente
detectável), ao passo que a proteína cai a $1/8$ e continua caindo à
medida que a proteína antiga se renova; a observação original de que a
proteína sumia ``enquanto o mRNA permanecia'' refletia sobretudo a
metade traducional da repressão.
\textbf{7.} $10^{6}/250 = 4000$ moléculas por óvulo.
\textbf{8.} $4000/550 \approx 7$ por célula ao eclodir. $4000/2^{n} < 1$
para $n \ge 12$ divisões desde o zigoto --- cerca de três divisões
depois da eclosão.
\textbf{9.} A polimerase copia o mensageiro-alvo em novo RNA de fita
dupla, cortado em \omterm{def:b3:rna-regulation:ncrna}{siRNAs} secundários: o gatilho é regenerado a partir do
próprio alvo, e por isso sobrevive à diluição e se espalha (um canal
passa RNA de fita dupla entre as células e para a linhagem germinativa).
Ele se desvanece porque a amplificação precisa ao mesmo tempo do alvo e
de um gatilho primário; uma vez ido o RNA injetado, a reserva secundária
é diluída geração após geração e a linhagem germinativa se reinicia.
\textbf{10.} $5000/2^{n} < 100$: $2^{n} > 50$, $n = 6$ dias (sobram $78$
cópias).
\textbf{11.} Células que não se dividem não diluem o \omterm{def:b3:rna-regulation:rnai}{RISC}; o limite é o
tempo de vida químico do \omterm{def:b3:rna-regulation:ncrna}{siRNA} (degradação por nucleases, retardada por
modificação química das fitas) e o vazamento lento do depósito endossomal
que abastece o citoplasma.
\textbf{12.} $2\times 10^{7}/\num{16384} \approx \num{1200}$
correspondências ao acaso. A genética decidiu: a perda de \emph{lin-14}
dá o fenótipo oposto ao da perda de \emph{lin-4}, os alelos de ganho de
função de \emph{lin-14} eram deleções exatamente dos sítios da 3$'$ UTR,
e os sete sítios eram conservados.
\textbf{13.} Nível $\propto$ meia-vida: subida de $50/10 = 5$ vezes.
\textbf{14.} Queda de $1/0.05 = 20$ vezes na síntese de ferritina.
\textbf{15.} A razão importação/armazenamento sobe $5\times 20 = 100$
vezes do estado rico ao pobre em ferro.
\textbf{16.} O agrupamento só se monta quando há ferro citosólico
disponível, de modo que a conformação da proteína é uma leitura direta da
concentração de ferro, sem receptor, hormônio ou transcrição --- como um
ribosswitch, um metabólito percebido pela molécula que ele controla.
\textbf{17.} A ferritina é traduzida independentemente do ferro: ferritina
sérica muito alta, ferro sérico e estoques normais (sem sobrecarga). No
cristalino, que não se renova, a ferritina se acumula por anos e
cristaliza --- uma catarata precoce.
\textbf{18.} O grampo é apenas um sítio de pouso; o efeito é aquilo que a
proteína ligada obstrui fisicamente --- os ribossomos que varrem a
extremidade 5$'$, uma nuclease na extremidade 3$'$.
\textbf{19.} Duas químicas cobrem duas condições: o agrupamento
ferro--enxofre da IRP1 também é destruído por oxidantes e por baixo
oxigênio, de modo que a IRP1 responde ao estado redox, ao passo que a
degradação da IRP2 dependente de ferro responde ao próprio ferro em toda
a faixa fisiológica de oxigênio; a redundância torna o regulão robusto, e
as duas relatam coisas diferentes.
\textbf{20.} $12\times 48\times 33\times 2 = \num{38016}$. $1 - (1 -
50/\num{38016})^{50} = 1 - 0.99868^{50} \approx 1 - e^{-0.066} =
0.064$: cerca de \qty{6}{\%}.
\textbf{21.} Isoformas idênticas ligam-se homofilicamente e se repelem: os
ramos de um mesmo neurônio se evitam e se espalham, ao passo que ramos de
neurônios diferentes, que quase não compartilham isoformas, podem se
cruzar --- uma identidade privada por célula.
\textbf{22.} $2^{n}$; $\num{1024}$ para $n = 10$. Menos na realidade: a
inclusão é decidida por fatores específicos de tipo celular e coordenada
entre os éxons, muitas combinações deslocam a fase de leitura e são
destruídas pelo decaimento mediado por códon sem sentido, e os éxons não
são independentes.
\textbf{23.} O \emph{Sxl} também controla a \omterm{def:b3:chromatin-epigenetics:x-inactivation}{compensação de dose}: nas
fêmeas a SXL bloqueia a tradução do \emph{msl-2}, de modo que os dois
cromossomos X não são hipertranscritos; sem SXL um embrião XX monta o
complexo MSL nos dois cromossomos X e dobra sua produção --- uma dose
letal de cada gene ligado ao X, antes mesmo de o sexo estar em questão.
\textbf{24.} O promotor precoce dispara apenas enquanto o sinal X:A está
presente; a SXL que ele produz força o splicing produtivo do transcrito
do promotor de manutenção, que é ativo nos dois sexos, de modo que, uma
vez existindo, a SXL continua produzindo a si mesma e o sinal deixa de ser
necessário --- a mesma lógica autossustentada de leitor-recruta-escritor
das marcas de \omterm{def:b3:chromatin-epigenetics:nucleosome}{cromatina}, aqui com uma proteína e um splicing.
\textbf{25.} Proteína LIN-14 reprimida $8$ vezes; meia-resposta do
interruptor do mensageiro \qty{10}{min}; um RNA não amplificado cai abaixo
de uma molécula por célula após $12$ divisões.
\end{solution}
